\section{A uniform non-alignment bound}\label{sec:rotation_proof}
We prove the main-text angle extension over the template and mass
box in~\eqref{eq:recovery_box} at $\alpha=.001$.
Keep the coefficient bounds~\eqref{eq:ind_coeff}. Let $f=.020$,
$B_0=\operatorname{diag}(\alpha,1)$, and
$\Delta B=B_\varphi-B_0$. For $|\varphi|\le f$ we have
$0\le\Delta B_{11}=-\Delta B_{22}\le f^2$ and
$|\Delta B_{12}|\le f$. Use $e$ and $E$ as in~\eqref{eq:ind_event},
with the rotated matrix. At fixed $e$, the rotated map differs
from the aligned map with $b_0=B_0\mu_z+e$ by
$-\gamma\Delta B(x-\mu_z)$.
The rotated map is nondecreasing in each output's own input
coordinate by the same covariance derivative as in S5, since
$B_{\varphi,jj}\le1$. Thus each face estimate extends across its
rectangle for the corresponding output coordinate.

\paragraph{Bounded common barrier.}
Let $W=[-.5,.3]\times[-.027,.51]$ for class 1.
Write $a_\pm,d_\pm,r_\pm,h_\pm$ for the endpoints of the
respective intervals in~\eqref{eq:recovery_box}.
For any $x=(s,t)\in W$, let $R=w_1/w_0$ and $H=w_2/w_0$.
The third-template half squared-distance gap is at least
\[
 Q_{\min}
 =\tfrac12(a_-^2+d_-^2-a_+-2d_+(.51))=.9002 .
\]
For $v\le.17$ this gives $H\le h_+e^{-Q_{\min}/.17}$,
so $w_2<.01$ and $\bar\mu_{v,1}>-.30$.
For $v\ge2.5$, split the range into $[2.5,4]$, $[4,8]$, and
$[8,\infty)$, each with endpoints $l,u$. On each piece,
\[
 R\ge R_* =r_-e^{-1/l},\qquad H\le H^*=h_+e^{-Q_{\min}/u}.
\]
Here $e^{-Q_{\min}/\infty}=1$. On all three pieces,
$H^*/(1+R_*+H^*)<.481$ and
$(R_*-a_+H^*)/(1+R_*+H^*)>-.30$.
Thus throughout $W$ the high and low ranges both satisfy
$w_2<.481$ and $\bar\mu_{v,1}>-.30$.

On the right face $s=.3$ at low scales, dropping $H$ yields
$\bar\mu_{v,1}\le L(.17)<.287729$.
For high scales split $t$ into $[-.027,.15]$ and $[.15,.51]$.
On a piece $[\ell,u]$ the gap is at most
\[
 Q^*=\tfrac12(a_+^2+d_+^2+.6a_+-2d_*\ell),
 \qquad
 d_*=\begin{cases}d_+&\ell<0,\\d_-&\ell\ge0.\end{cases}
\]
Both upper bounds are positive, so
$H\ge h_-e^{-Q^*/2.5}$ and $R\le r_+$.
Substituting these bounds into $(R-aH)/(1+R+H)$ gives
$\bar\mu_{v,1}<.29300$ and $\bar\mu_{v,1}<.25984$ on the
two pieces, respectively. Denote the respective upper bounds by $M$. With $U=\max(|\ell|,|u|)$
on each high-scale piece and $U=.51$ at low scales, the
right-face drift is at most
\[
 \gamma[.00095+.999(M-.3)+.7f^2+fU].
\]
The three upper bounds (low, first high, second high) are below
$-.00020$, $-.00069$, and $-.0071$, respectively.
The left-face drift is at least
$\gamma[.0015-.00025+.999(.5-.30)-f(.51)]>.047$.
The bottom-face drift is at least
$\gamma[-.005+1.999(.027)-1.5f-.027f^2]>.0047$.
The top-face drift is at most
$\gamma[.005-.51+1.002(2.02(.481)-.51)+1.5f+.51f^2]<-.0030$.
For class 1 the initial coordinates on $E$ satisfy
$.00075\le b_1\le.00165$ and $|b_2|\le.025$.
The four margins and the initialization prove that all
restricted trajectories stay in $W$.

\paragraph{Entry under the preset.}
The class-0 and class-1 backprojections retain their original
decision thresholds. For class 2,
$\|\Delta B\mu_2\|_\infty\le(f+f^2)2.02$ implies
$|b_1|<.043$ and $1.933<b_2<2.067$. The classifier therefore
identifies class 2 and its backprojection belongs to
$G_{2,\varphi}=[-1.12,.05]\times[1.73,2.27]$.
Class 0 starts in $G_0$. Both classes use only $v_L$.

For class 1, write $s_-=-.01$, $\xi=.53$, $s_+=1.1$,
$T=.06$, and $v_-=.69,v_+=.71$.
The channel and prepared strip share $-.027\le t\le T$ and have
horizontal intervals $[s_-,\xi]$ and $[\xi,s_+]$, respectively.
On their union,
\begin{align*}
 R_{\rm lo}&=r_-\exp((s_--.5)/v_-),\\
 H_{\rm hi}&=h_+\exp[-(a_-^2+d_-^2+2a_+s_--2d_+T)/(2v_+)]
\end{align*}
bound the two odds from below and above, respectively. Thus
$w_2\le H_{\rm hi}/(1+R_{\rm lo}+H_{\rm hi})<.04508$,
and $\partial_s\bar\mu_{v,1}\le(1/4+a_+(a_++2)w_2)/v_-<.564$.
At $s=\xi$, the lower bound on $R$ becomes
$r_-e^{(\xi-.5)/v_+}$ and the upper bound on $H$ is obtained
by replacing $2a_+s_-$ by $2a_-\xi$. These give a surplus
$\bar\mu_{v,1}(\xi,t)-\xi>.01563$.
The horizontal and vertical rotation corrections on the union
are bounded by $\gamma(1.01 f^2+Tf)<.000402$ and
$\gamma(1.01f+Tf^2)=.005056$, respectively.
Using the $w_2$ bound, the top-face displacement is below
$-.0009$, and the bottom-face displacement exceeds $.007$.
The aligned data force is nonnegative up to $\xi$.
The lower comparison from $\underline s_0=.00075$ is
\[
 \underline s_{k+1}
 =\underline s_k+\gamma(.999)
       [.01563+.436(\xi-\underline s_k)]-.000402 .
\]
This yields $\underline s_{25}>.5306$. The same positive drift
on $s=\xi$ keeps the left face of the prepared strip invariant.
Since $\bar\mu_{v,1}\le1$, the overshoot is below $.649$ and
the right face $s=1.1$ is inward even after adding $.000402$.
The prepared strip lies in $G_1$, so the 25-call switch is valid
for every intermediate history in $[.69,.71]$.
The numerical face and comparison bounds above follow by
outward-rounded evaluation of the displayed scalar expressions.

\paragraph{Actual terminal dynamics and risk.}
The enlargement to $G_{2,\varphi}$ preserves the gap bound $.06$:
on that rectangle the gap against $\mu_0$ is at least
$2(1.98)(1.73)-2(1.02)(.05)-1.02^2-2.02^2>1.62$,
and the gap against $\mu_1$ adds $1-2s\ge.9$.
The other rectangles retain their original gaps.
All terminal rectangles have $\|x-\mu_z\|_\infty\le1.1$.
The rotation correction is bounded there by
$\gamma(f+f^2)1.1=.00561$, which fits the aligned terminal
face margin $\gamma(.999\cdot.09-.005-1.002\cdot12e^{-30})>.021$.
Hence the rotated terminal map preserves all three rectangles.

Use the actual center
$m_\varphi=(B_\varphi+c_LI)^{-1}(b+c_L\mu_z)$.
The residual recursion uses $I-\gamma(B_\varphi+c_LI)$, whose
norm is at most $.751$, and forcing
$\gamma c_L(\bar\mu_{v_L}-\mu_z)$ of norm at most
$12\gamma c_Le^{-.03/v_L}$. Summing the recursion yields
$O(q^n+e^{-.03/v_L})$ with $q=.751$.
The actual terminal map also contracts on the rectangles by
$\|C_{v_L}\|\le36e^{-.03/v_L}$ and $v_L\le.001$.

Conditional on the class, $e$ has covariance
$\tau^2B_\varphi^2+\nu^2B_\varphi\preceq S^2I$,
so $\Prob(E^c)\le4\exp[-(.00025)^2/(2S^2)]$ uniformly in $\varphi$.
The affine-center risk is unchanged because the two sensing
eigenvalues and isotropic covariances are preserved.
On all observations and at every admissible scale the linear
part has norm at most $.751$ and the forcing is bounded by
a constant times $1+\|b\|$. Iterates and centers have
budget-uniform fourth moments for $S\le1$.
Cauchy--Schwarz controls the squared residual on the
exceptional event. Taking $\lambda=1$, $v_L=S$, and the
terminal budget of~\eqref{eq:fixed_geom} yields $O(S^2)$.
On $E$ for class 1 the restricted output still satisfies $s\le.3$,
so the same directional lower bound gives the limiting $.147$.

\paragraph{Reproduction.}
From the artifact root, run
\path{python experiments/code/rotation_certificate.py}.
The only nonstandard dependency is \texttt{mpmath} (1.3.0).
The expected output is \texttt{PASS: 57 strict scalar bounds at
45 digits}. The optional flag \path{--out bounds.json} records
both endpoints of each comparison. The check uses outward-rounded
arithmetic on the scalar reductions above, without trajectory
sampling. The argument relies on the stated fixed gain and
gives the risk rate in~\eqref{eq:fixed_geom}; the calibrated
oracle excess estimate has its own hypotheses.
